\documentclass[aspectratio=169,10pt]{beamer}

% Shared Beamer setup for Fall 2026 course slides.
\mode<presentation>{
  \usetheme{Madrid}
  \usecolortheme{default}
}

\definecolor{CalPolyGreen}{HTML}{154734}
\definecolor{CalPolyGold}{HTML}{BD8B13}
\definecolor{DeckInk}{HTML}{1F2933}
\definecolor{DeckMuted}{HTML}{5F6B7A}

\setbeamercolor{structure}{fg=CalPolyGreen}
\setbeamercolor{palette primary}{bg=CalPolyGreen,fg=white}
\setbeamercolor{palette secondary}{bg=DeckInk,fg=white}
\setbeamercolor{palette tertiary}{bg=CalPolyGold,fg=DeckInk}
\setbeamercolor{frametitle}{bg=CalPolyGreen,fg=white}
\setbeamercolor{title}{fg=CalPolyGreen}
\setbeamercolor{normal text}{fg=DeckInk,bg=white}
\setbeamercolor{block title}{bg=CalPolyGreen,fg=white}
\setbeamercolor{block body}{bg=black!3,fg=DeckInk}

\setbeamertemplate{navigation symbols}{}
\setbeamertemplate{footline}[frame number]
\setbeamertemplate{itemize item}{\color{CalPolyGold}$\blacktriangleright$}
\setbeamertemplate{itemize subitem}{\color{CalPolyGreen}$\bullet$}

\newcommand{\CourseNumber}{Course}
\newcommand{\CourseTitle}{Course Title}
\newcommand{\LectureNumber}{0}
\newcommand{\LectureTitle}{Lecture Title}
\newcommand{\LectureDate}{Date}
\newcommand{\CourseUnit}{Unit}

\newcommand{\coursenumber}[1]{\renewcommand{\CourseNumber}{#1}}
\newcommand{\coursetitle}[1]{\renewcommand{\CourseTitle}{#1}}
\newcommand{\lecturenumber}[1]{\renewcommand{\LectureNumber}{#1}}
\newcommand{\lecturetitle}[1]{\renewcommand{\LectureTitle}{#1}}
\newcommand{\lecturedate}[1]{\renewcommand{\LectureDate}{#1}}
\newcommand{\courseunit}[1]{\renewcommand{\CourseUnit}{#1}}

\author{Kirk Duran}
\institute{California Polytechnic State University, San Luis Obispo}
\date{\LectureDate}

\newcommand{\makedecktitle}{
  \begin{frame}[plain]
    \vspace{0.25in}
    {\usebeamerfont{title}\color{CalPolyGreen}\LectureTitle\par}
    \vspace{0.18in}
    {\large \CourseNumber{}: \CourseTitle\par}
    \vspace{0.08in}
    {\large Lecture \LectureNumber{} -- \LectureDate\par}
    \vfill
    {\small Kirk Duran \textbar{} Cal Poly, San Luis Obispo\par}
  \end{frame}
}

\newcommand{\placeholdernote}{
  \begin{block}{Placeholder}
    Replace these bullets with the final lecture material, examples, and in-class activities.
  \end{block}
}

\AtBeginSection[]{
  \begin{frame}{Roadmap}
    \tableofcontents[currentsection]
  \end{frame}
}


\usepackage{tikz}

\graphicspath{{../assets/lecture-06/}}

% If an image file is not yet available, skip the visual slot.
\newcommand{\lectureimage}[2][1.75in]{%
  \IfFileExists{../assets/lecture-06/#2}{%
    \includegraphics[width=\linewidth,height=#1,keepaspectratio]{#2}%
  }{%
  }%
}

% Temporary visible placeholder used while lecture assets are being assembled.
% Replace each placeholder with \lectureimage[...]{} once the corresponding asset exists.
\newcommand{\lectureplaceholder}[2][1.75in]{%
  \begin{center}
    \fbox{%
      \parbox[c][#1][c]{0.90\linewidth}{%
        \centering\small\textcolor{DeckMuted}{\textbf{Image placeholder}\\[0.08in]#2}%
      }%
    }
  \end{center}%
}

\setbeamersize{text margin left=0.42in,text margin right=0.42in}

\coursenumber{CS 4880}
\coursetitle{Artificial Intelligence}
\lecturenumber{6}
\lecturetitle{Heuristic Search}
\lecturedate{September 4, 2026}
\courseunit{Search and Problem Solving}

\begin{document}

\makedecktitle

\begin{frame}[shrink=20]{Today}
  \begin{itemize}
    \item Why uninformed search wastes effort when the state space is large.
    \item How a heuristic function estimates remaining cost to the goal.
    \item How Greedy Best-First Search uses only heuristic guidance.
    \item How A* combines path cost and heuristic guidance.
    \item Why admissibility and consistency matter for A*.
    \item How Weighted A* and related variants trade solution quality for search effort.
  \end{itemize}

  \begin{keyidea}
    Informed search uses problem-specific estimates to decide which frontier states appear most promising.
  \end{keyidea}
\end{frame}

\sectionframe{Part 1: From Uninformed Search to Informed Search}

\begin{frame}[shrink=20]{A Human Would Look Toward the Goal}
  \begin{columns}[T,onlytextwidth]
    \begin{column}{0.53\textwidth}
      Suppose we are navigating a city and trying to reach a visible landmark.

      \begin{itemize}
        \item BFS can systematically enumerate nearby intersections.
        \item Dijkstra can account for travel cost along roads.
        \item A person can also use the visible landmark to estimate which direction is promising.
      \end{itemize}

      \begin{block}{Useful extra information}
        The landmark provides an estimate of how far we still appear to be from the goal.
      \end{block}
    \end{column}

    \begin{column}{0.41\textwidth}
      \lectureimage[2.55in]{slide6.png}
    \end{column}
  \end{columns}

  \begin{keyidea}
    Search can exploit information that is not encoded by graph connectivity alone.
  \end{keyidea}
\end{frame}

\begin{frame}[shrink=20]{The Search Problem Solver Problem}
  \begin{columns}[T,onlytextwidth]
    \begin{column}{0.55\textwidth}
      Uninformed search uses only the structure of the problem space.

      \begin{itemize}
        \item BFS, DFS, and iterative deepening do not use problem-specific estimates of progress.
        \item In deep or highly branching spaces, large amounts of work may be spent in irrelevant regions.
        \item Even simple puzzle domains can contain enormous state spaces.
      \end{itemize}

      \begin{block}{Example}
        The 15-puzzle has more than $20$ trillion reachable states.
      \end{block}
    \end{column}

    \begin{column}{0.39\textwidth}
      \lectureimage[2.48in]{slide7.png}
    \end{column}
  \end{columns}

  \begin{keyidea}
    Systematic search is not enough by itself; large problems require a way to prioritize promising states.
  \end{keyidea}
\end{frame}

\begin{frame}[shrink=20]{Cognitive Psychology Meets AI}
  \begin{columns}[T,onlytextwidth]
    \begin{column}{0.53\textwidth}
      The source lecture motivates heuristics using observations from human problem solving.

      \begin{itemize}
        \item People estimate how far they appear to be from a goal.
        \item People often avoid paths that appear to move in the wrong direction.
        \item Human problem solving is therefore not purely blind enumeration.
      \end{itemize}

      \begin{block}{AI question}
        How can we encode ``how close does this state appear to be to the goal?''
      \end{block}
    \end{column}

    \begin{column}{0.41\textwidth}
      \lectureimage[2.55in]{slide8.png}
    \end{column}
  \end{columns}

  \begin{keyidea}
    A heuristic converts domain knowledge into a numerical estimate that can guide search order.
  \end{keyidea}
\end{frame}

\sectionframe{Part 2: Heuristic Functions --- Estimating Cost-to-Go}

\begin{frame}[shrink=20]{Heuristic Function}
  \begin{cpdefinition}
    A \textbf{heuristic function} $h(n)$ estimates the cost of a cheapest path from state $n$ to a goal state.
  \end{cpdefinition}

  \begin{columns}[T,onlytextwidth]
    \begin{column}{0.52\textwidth}
      \begin{itemize}
        \item Smaller $h(n)$ means the state appears closer to the goal.
        \item Heuristics are typically designed for a particular problem domain.
        \item A good heuristic can reduce the number of states that must be expanded.
      \end{itemize}

      \begin{block}{Interpretation}
        $h(n)$ predicts future cost; it does not measure the cost already incurred.
      \end{block}
    \end{column}

    \begin{column}{0.42\textwidth}
      \lectureimage[2.42in]{slide10.png}
    \end{column}
  \end{columns}

  \begin{keyidea}
    The purpose of a heuristic is not to solve the problem directly, but to rank frontier states more intelligently.
  \end{keyidea}
\end{frame}

\begin{frame}[shrink=20]{Common Heuristics: Grids and Sliding-Tile Puzzles}
  \begin{columns}[T,onlytextwidth]
    \begin{column}{0.48\textwidth}
      \begin{block}{Grid navigation}
        \begin{itemize}
          \item Euclidean distance
          \item Manhattan distance
        \end{itemize}
      \end{block}

      For states $(x_1,y_1)$ and $(x_2,y_2)$:
      \[
        h_{\mathrm{Manhattan}}=|x_1-x_2|+|y_1-y_2|.
      \]
    \end{column}

    \begin{column}{0.48\textwidth}
      \begin{block}{8-puzzle}
        \begin{itemize}
          \item number of misplaced tiles;
          \item sum of Manhattan distances from tiles to their goal positions.
        \end{itemize}
      \end{block}

      \lectureimage[1.60in]{slide11.png}
    \end{column}
  \end{columns}

  \begin{keyidea}
    Different heuristics summarize different approximations of the remaining work.
  \end{keyidea}
\end{frame}

\begin{frame}[shrink=20]{Common Heuristics Across AI Domains}
  \begin{columns}[T,onlytextwidth]
    \begin{column}{0.30\textwidth}
      \begin{block}{Traveling Salesperson}
        Minimum spanning tree over unvisited cities.
      \end{block}
    \end{column}

    \begin{column}{0.32\textwidth}
      \begin{block}{Robot Motion Planning}
        Straight-line distance to the target while ignoring obstacles.
      \end{block}
    \end{column}

    \begin{column}{0.32\textwidth}
      \begin{block}{Planning / STRIPS}
        Number of goal conditions or preconditions that remain unsatisfied.
      \end{block}
    \end{column}
  \end{columns}

  \vspace{0.12in}
  \lectureimage[1.35in]{slide12.png}

  \begin{keyidea}
    Heuristic design is problem-specific: the estimate should capture structure that correlates with true remaining cost.
  \end{keyidea}
\end{frame}

\begin{frame}[shrink=20]{Making Heuristic Functions}
  \begin{columns}[T,onlytextwidth]
    \begin{column}{0.54\textwidth}
      \begin{block}{1. Solve a relaxed problem}
        Remove constraints so that the remaining problem is easier to solve.
        \begin{itemize}
          \item ignore obstacles;
          \item ignore interactions;
          \item ignore selected costs or restrictions.
        \end{itemize}
      \end{block}

      \begin{block}{2. Use domain knowledge}
        Encode expert knowledge about distance, effort, or utility.
      \end{block}

      \begin{block}{3. Learn the heuristic}
        Train a model to approximate cost-to-go from solved instances.
      \end{block}
    \end{column}

    \begin{column}{0.40\textwidth}
      \lectureimage[2.65in]{slide13.png}
    \end{column}
  \end{columns}
\end{frame}

\sectionframe{Part 3: Greedy Best-First Search --- Follow the Heuristic}

\begin{frame}[shrink=20]{Greedy Best-First Search}
  \begin{cpdefinition}
    \textbf{Greedy Best-First Search (GBFS)} expands the frontier node with the smallest heuristic value $h(n)$.
  \end{cpdefinition}

  \[
    \mathrm{priority}(n)=h(n)
  \]

  \begin{itemize}
    \item GBFS attempts to reach the goal quickly by following the state that appears closest.
    \item It repeatedly removes the minimum-$h$ state, expands it, and inserts its successors.
    \item It ignores accumulated path cost $g(n)$ when choosing what to expand next.
  \end{itemize}

  \begin{keyidea}
    GBFS is greedy because it optimizes only its estimate of the remaining path, not the total solution cost.
  \end{keyidea}
\end{frame}

\begin{frame}[shrink=20]{Greedy Best-First Search Pseudocode}
  \begin{block}{Algorithm}
  \ttfamily\small
  frontier $\leftarrow$ PriorityQueue(order by $h$)\\
  frontier.insert(initial\_state, priority=$h(initial\_state)$)\\
  came\_from[initial\_state] $\leftarrow$ None\\[0.06in]
  while frontier is not empty:\\
  \quad state $\leftarrow$ frontier.remove\_min()\\[0.03in]
  \quad if goal\_test(state):\\
  \qquad return reconstruct\_path(came\_from, state)\\[0.03in]
  \quad for child in SUCCESSORS(state):\\
  \qquad if child has not been explored:\\
  \qquad\quad came\_from[child] $\leftarrow$ state\\
  \qquad\quad frontier.insert(child, priority=$h(child)$)\\[0.03in]
  return failure
  \end{block}

  \begin{keyidea}
    The priority queue contains only heuristic estimates; edge costs do not affect expansion order.
  \end{keyidea}
\end{frame}

\begin{frame}[shrink=20]{GBFS Worked Example: From \texttt{A} to \texttt{Z}}
  \begin{columns}[T,onlytextwidth]
    \begin{column}{0.52\textwidth}
      Source heuristic values include:
      \[
        h(A)=14,\quad h(C)=11,\quad h(E)=4,\quad h(Z)=0.
      \]

      \begin{block}{Expansion sequence}
        \centering
        $A:14$\\
        $\downarrow$\\
        $C:11$\\
        $\downarrow$\\
        $E:4$\\
        $\downarrow$\\
        $Z:0$
      \end{block}

      GBFS chooses the smallest available $h$ value at every step.
    \end{column}

    \begin{column}{0.42\textwidth}
      \lectureimage[2.55in]{slide18.png}
    \end{column}
  \end{columns}

  \begin{keyidea}
    The route looks promising according to $h$, even though GBFS does not compare the path costs used to reach those states.
  \end{keyidea}
\end{frame}

\begin{frame}[shrink=20]{Uninformative Heuristic Regions}
  \begin{columns}[T,onlytextwidth]
    \begin{column}{0.55\textwidth}
      An \textbf{uninformative region} occurs when many distinct states receive similar or identical heuristic values.

      \begin{itemize}
        \item The heuristic cannot clearly distinguish which states are more promising.
        \item Greedy Best-First Search or A* may expand many states unnecessarily.
        \item Search can begin to resemble uninformed exploration.
        \item Large flat regions can produce substantial slowdowns.
      \end{itemize}
    \end{column}

    \begin{column}{0.39\textwidth}
      \lectureimage[2.55in]{slide14.png}
    \end{column}
  \end{columns}

  \begin{keyidea}
    A heuristic is useful only to the extent that its values meaningfully rank competing frontier states.
  \end{keyidea}
\end{frame}

\sectionframe{Part 4: A* Search --- Balance Cost-So-Far and Cost-to-Go}

\begin{frame}[shrink=20]{A* Combines Two Signals}
  \begin{columns}[T,onlytextwidth]
    \begin{column}{0.52\textwidth}
      A* evaluates each frontier node using:
      \[
        f(n)=g(n)+h(n).
      \]

      \begin{itemize}
        \item $g(n)$: exact path cost from the start to $n$;
        \item $h(n)$: estimated cost from $n$ to a goal;
        \item $f(n)$: estimated total solution cost through $n$.
      \end{itemize}
    \end{column}

    \begin{column}{0.42\textwidth}
      \begin{block}{Special cases}
        \begin{itemize}
          \item If $h(n)=0$, A* behaves like Dijkstra / UCS.
          \item If $g(n)$ is ignored, the policy becomes GBFS.
        \end{itemize}
      \end{block}

      \lectureimage[1.55in]{slide21.png}
    \end{column}
  \end{columns}

  \begin{keyidea}
    A* prefers nodes that are inexpensive so far \emph{and} appear inexpensive to finish.
  \end{keyidea}
\end{frame}

\begin{frame}[shrink=20]{A* Search Pseudocode}
  \begin{block}{Algorithm}
  \ttfamily\small
  frontier $\leftarrow$ PriorityQueue(order by $f$)\\
  frontier.insert(start, priority=$h(start)$)\\
  cost\_so\_far[start] $\leftarrow 0$\\
  came\_from[start] $\leftarrow$ None\\[0.06in]
  while frontier is not empty:\\
  \quad state $\leftarrow$ frontier.remove\_min()\\
  \quad if goal\_test(state): return reconstruct\_path(...)\\[0.03in]
  \quad for child in SUCCESSORS(state):\\
  \qquad new\_cost $\leftarrow g(state)+c(state,child)$\\
  \qquad if child is new or new\_cost $< g(child)$:\\
  \qquad\quad $g(child)\leftarrow$new\_cost\\
  \qquad\quad $f(child)\leftarrow g(child)+h(child)$\\
  \qquad\quad frontier.insert(child, priority=$f(child)$)
  \end{block}
\end{frame}

\begin{frame}[shrink=20]{A* Worked Example: Initial Expansions}
  \begin{columns}[T,onlytextwidth]
    \begin{column}{0.55\textwidth}
      Start at $A$:
      \[
        f(A)=0+14=14.
      \]

      Expand $A$:
      \[
        f(C)=3+11=14,
      \]
      \[
        f(B)=4+12=16.
      \]

      Frontier:
      \[
        [C:14,\;B:16].
      \]

      Expand $C$:
      \[
        f(D)=10+6=16,\qquad f(E)=13+4=17.
      \]
    \end{column}

    \begin{column}{0.39\textwidth}
      \lectureplaceholder[2.55in]{A--Z graph with edge costs and heuristic values used by the A* example.}
    \end{column}
  \end{columns}
\end{frame}

\begin{frame}[shrink=20]{A* Worked Example: Updating the Best Route}
  Continuing the source example:

  \begin{itemize}
    \item After expanding $B$, the frontier contains $D:16$, $E:17$, and $F:20$.
    \item Expanding $D$ discovers a cheaper route to $E$:
      \[
        g(E)=12,\qquad f(E)=12+4=16.
      \]
    \item Expanding $E:16$ generates:
      \[
        f(Z)=17+0=17.
      \]
    \item When $Z:17$ is selected, the search terminates with solution cost $17$.
  \end{itemize}

  \begin{keyidea}
    Unlike GBFS, A* can revise a promising-looking state when a cheaper path to that state is discovered.
  \end{keyidea}
\end{frame}

\begin{frame}[shrink=20]{Admissible Heuristics}
  \begin{cpdefinition}
    A heuristic is \textbf{admissible} if it never overestimates the true minimum remaining cost.
  \end{cpdefinition}

  \[
    h(n)\le h^*(n)
  \]

  where $h^*(n)$ is the true minimum cost from $n$ to a goal.

  \begin{itemize}
    \item An admissible heuristic is optimistic: it may underestimate, but it does not claim the goal is farther away than it really is.
    \item The closer $h(n)$ is to $h^*(n)$ while remaining admissible, the more strongly it can guide search.
  \end{itemize}

  \begin{keyidea}
    Admissibility connects heuristic quality to A*'s minimum-cost guarantee.
  \end{keyidea}
\end{frame}

\begin{frame}[shrink=20]{Consistent Heuristics}
  \begin{cpdefinition}
    A heuristic is \textbf{consistent} (monotone) if for every transition from $n$ to successor $n'$:
  \end{cpdefinition}

  \[
    h(n)\le c(n,n')+h(n').
  \]

  \begin{itemize}
    \item The estimated remaining cost cannot decrease by more than the cost of the step just taken.
    \item Consistency is a triangle-inequality condition on heuristic estimates.
    \item In graph search, consistency supports the standard rule that a state need not be reopened after it has been finalized under the usual A* formulation.
  \end{itemize}

  \begin{keyidea}
    Consistency makes heuristic estimates behave coherently along paths, not merely at isolated states.
  \end{keyidea}
\end{frame}

\begin{frame}[shrink=20]{What A* Pays for Its Guidance}
  \begin{columns}[T,onlytextwidth]
    \begin{column}{0.47\textwidth}
      \begin{block}{Search costs}
        \begin{itemize}
          \item can store a very large open and closed set;
          \item high branching factors still create many candidates;
          \item enormous state spaces can remain impractical;
          \item priority-queue maintenance adds computation.
        \end{itemize}
      \end{block}
    \end{column}

    \begin{column}{0.47\textwidth}
      \begin{block}{Heuristic dependence}
        \begin{itemize}
          \item weak heuristics approach uninformed search;
          \item heuristic evaluation itself may be expensive;
          \item exact optimal search may be too slow for real-time decisions.
        \end{itemize}
      \end{block}
    \end{column}
  \end{columns}

  \begin{keyidea}
    A* improves search order, but it does not eliminate the combinatorial growth of difficult state spaces.
  \end{keyidea}
\end{frame}

\sectionframe{Part 5: Variants --- Trading Optimality, Memory, and Exploration}

\begin{frame}[shrink=20]{Related A* Variants and Tradeoffs}
  \begin{columns}[T,onlytextwidth]
    \begin{column}{0.47\textwidth}
      \begin{block}{Bidirectional A*}
        Search simultaneously from the start and goal to reduce the effective search depth.
      \end{block}

      \begin{block}{Iterative Deepening A* (IDA*)}
        Use heuristic cost thresholds with depth-first traversal to reduce memory requirements.
      \end{block}
    \end{column}

    \begin{column}{0.47\textwidth}
      \begin{block}{Beam Search}
        Keep only a limited number of promising frontier states, sacrificing completeness or optimality for memory and speed.
      \end{block}

      \begin{block}{Weighted A*}
        Increase the influence of the heuristic to search more aggressively toward the goal.
      \end{block}
    \end{column}
  \end{columns}
\end{frame}

\begin{frame}[shrink=20]{Weighted A*}
  \begin{cpdefinition}
    \textbf{Weighted A*} modifies the A* evaluation function by multiplying the heuristic by a weight $w>1$.
  \end{cpdefinition}

  \[
    f_w(n)=g(n)+w\,h(n).
  \]

  \begin{columns}[T,onlytextwidth]
    \begin{column}{0.47\textwidth}
      \begin{block}{Weight bias}
        \begin{itemize}
          \item $w=1$: ordinary A*;
          \item larger $w$: stronger preference for states that appear close to the goal;
          \item very large $w$: behavior approaches GBFS.
        \end{itemize}
      \end{block}
    \end{column}

    \begin{column}{0.47\textwidth}
      \begin{block}{Tradeoff}
        \begin{itemize}
          \item often expands fewer states;
          \item may return a suboptimal solution;
          \item can provide bounded suboptimality under appropriate heuristic assumptions.
        \end{itemize}
      \end{block}
    \end{column}
  \end{columns}
\end{frame}

\begin{frame}[shrink=20]{Weighted A* in Uninformative Heuristic Regions}
  \begin{columns}[T,onlytextwidth]
    \begin{column}{0.54\textwidth}
      Increasing $w$ makes the search behave more like GBFS.

      \begin{itemize}
        \item The influence of $g(n)$ becomes smaller relative to $w\,h(n)$.
        \item In a flat heuristic region, many states still receive nearly indistinguishable priorities.
        \item The search can lose direction despite using a larger heuristic weight.
        \item More aggressive weighting also increases the expected loss of optimality.
      \end{itemize}
    \end{column}

    \begin{column}{0.40\textwidth}
      \lectureimage[2.55in]{slide31.png}
    \end{column}
  \end{columns}

  \begin{keyidea}
    Increasing heuristic influence helps only when the heuristic contains useful information to amplify.
  \end{keyidea}
\end{frame}

\begin{frame}[shrink=20]{Exploration vs. Exploitation in Greedy Search}
  \begin{columns}[T,onlytextwidth]
    \begin{column}{0.58\textwidth}
      The source deck introduces mechanisms that deliberately add exploration to greedy search:

      \begin{itemize}
        \item \textbf{Randomized GBFS}: occasionally choose among several best-ranked states.
        \item \textbf{$\varepsilon$-greedy GBFS}: with probability $\varepsilon$, choose a random open state; otherwise choose minimum $h$.
        \item \textbf{Multiple queues}: maintain frontier views organized by different metadata or search criteria.
        \item \textbf{Type-based GBFS}: group states into domain-defined types and sample across groups.
        \item \textbf{Diverse best-first search}: favor dissimilar states among highly ranked candidates.
      \end{itemize}
    \end{column}

    \begin{column}{0.36\textwidth}
      \lectureimage[2.55in]{slide32.png}
    \end{column}
  \end{columns}
\end{frame}

\begin{frame}[shrink=20]{Exploration Mechanisms in Weighted A*}
  \begin{columns}[T,onlytextwidth]
    \begin{column}{0.30\textwidth}
      \begin{block}{Randomized focal selection}
        Break ties or select randomly within a focal list to encourage diversity among similarly promising states.
      \end{block}
    \end{column}

    \begin{column}{0.32\textwidth}
      \begin{block}{$\varepsilon$-greedy focal WA*}
        With probability $\varepsilon$, select a random focal state; otherwise optimize a secondary heuristic.
      \end{block}
    \end{column}

    \begin{column}{0.32\textwidth}
      \begin{block}{Type-based Weighted A*}
        Group focal states by domain-specific type and favor under-explored structural categories.
      \end{block}
    \end{column}
  \end{columns}

  \begin{keyidea}
    Modern variants can combine bounded-cost search with explicit mechanisms for diversity and exploration.
  \end{keyidea}
\end{frame}

\begin{frame}[shrink=20]{Search Policies at a Glance}
  \begin{center}
  \renewcommand{\arraystretch}{1.35}
  \begin{tabular}{l|c|l}
    \textbf{Algorithm} & \textbf{Priority} & \textbf{Main behavior} \\\hline
    Dijkstra / UCS & $g(n)$ & cheapest path discovered so far \\
    GBFS & $h(n)$ & state that appears closest to goal \\
    A* & $g(n)+h(n)$ & balance cost-so-far and estimated cost-to-go \\
    Weighted A* & $g(n)+w h(n)$ & favor heuristic guidance more aggressively
  \end{tabular}
  \end{center}

  \begin{keyidea}
    The evaluation function determines what the search algorithm considers ``promising.''
  \end{keyidea}
\end{frame}

\begin{frame}[shrink=20]{Takeaways}
  \begin{itemize}
    \item Uninformed algorithms differ mainly in how they order states using depth or accumulated cost.
    \item A heuristic $h(n)$ estimates the remaining cost to a goal and injects domain knowledge into search.
    \item GBFS uses only $h(n)$, which can be fast but can also be misled.
    \item A* uses $f(n)=g(n)+h(n)$ to combine actual cost and estimated remaining cost.
    \item Admissibility and consistency formalize when heuristic guidance is safe for optimal graph search.
    \item Weighted and exploratory variants deliberately trade exact optimality, memory, or determinism for practical search performance.
  \end{itemize}

  \begin{keyidea}
    Better search comes from matching the frontier priority rule to the information and guarantees the problem actually requires.
  \end{keyidea}
\end{frame}

\end{document}
